
\begin{frame}[fragile]{Valeurs nulles}

Une \alert{valeur nulle} est une
\alert{valeur manquante}. Ne pas confondre avec une chaîne de caractères  ``null''  ou 
vide (``'')

\vfill

Voici une table des logements avec quelques valeurs à \sqlkw{null}

\vfill
\begin{tabular}{c|c|c|c|c}
\textbf{code} & \textbf{nom} & \textbf{capacité} & \textbf{type} & \textbf{lieu}\\ \hline 
ca & Causses & 45 & Auberge & Cévennes\\ \hline 
ge & Génépi & 134 & Hôtel & \\ \hline 
pi & U Pinzutu &  & Gîte & Corse\\ \hline 
ta & Tabriz & 34 & Hôtel & Bretagne\\ \hline 
\end{tabular}

\vfill

La présence de valeurs nulles \alert{fausse} le résultat attendu des requêtes.

\end{frame}


\begin{frame}[fragile]{Comparaisons avec valeurs nulles}

On ne sait ni comparer une valeur nulle, ni appliquer une fonction.

\vfill


\begin{minted}[baselinestretch=.9,
               fontsize=\small,
               framesep=2mm]{sql}
select * from Logement where lieu = ''
\end{minted}

Génépi (pas de lieu) n'est pas trouvé.

\vfill

\begin{minted}[baselinestretch=.9,
               fontsize=\small,
               framesep=2mm]{sql}
select * from Logement where lieu != ''
\end{minted}
Génépi (pas de lieu) n'est pas trouvé non plus!
\vfill

\begin{blocgauche}
Une comparaison avec \sqlkw{null} ne donne ni vrai, ni faux,
mais une troisième valeur de vérité, \sqlkw{unknown}. 
\end{blocgauche}
\end{frame}


%%%%%%%%
\begin{frame}[fragile]{Calculs avec valeurs à null}

Tout calcul appliqué à une valeur à \sqlkw{null} renvoie \sqlkw{null}
\vfill


\begin{minted}[baselinestretch=.9,
               fontsize=\small,
               framesep=2mm]{sql}
select nom, capacité * 2 as 'doubleCapacité'
from Logement
\end{minted}

\vfill
\begin{tabular}{c|c}
\textbf{nom} & \textbf{doubleCapacité}\\ \hline 
Causses & 90\\ \hline 
Génépi & 268\\ \hline 
U Pinzutu & \\ \hline 
Tabriz & 68\\ \hline 
\end{tabular}

\end{frame}

%%%%%%%%%%%%%%%%%%%%
\begin{frame}[fragile]{Le test \sqlkw{is null}}

Seule approche correcte: il faut tester explicitement l'absence de valeur avec \texttt{is null}.

\vfill
En SQL: 

\begin{minted}[baselinestretch=.9,
               fontsize=\small,
               framesep=2mm]{sql}
select *
from Logement
where capacité is null
\end{minted}

\vfill

\alert{Attention}\,: le test \texttt{capacité = null} \alert{ne marche pas}.

\vfill

\alert{Conclusion}\,: éviter autant que possible les valeurs nulles
en les interdisant (dans le schéma).
\end{frame}

%%%%%%%%%%%%%%%%%%%%
\begin{frame}[fragile]{La jointure externe}

On veut la liste des logements avec leurs activités.
\vfill


\begin{minted}[baselinestretch=.9,
               fontsize=\small,
               framesep=2mm]{sql}
select nom, codeActivité
from Logement, Activité
where code = codeLogement
\end{minted}

\vfill
\begin{footnotesize}
\begin{tabular}{c|c}
\textbf{nom} & \textbf{codeActivité}\\ \hline 
Causses & Randonnée\\ \hline 
Génépi & Piscine\\ \hline 
Génépi & Ski\\ \hline 
U Pinzutu & Plongée\\ \hline 
U Pinzutu & Voile\\ \hline 
\end{tabular}
\end{footnotesize}
\vfill

\begin{blocgauche}
\alert{Il manque l'hôtel Tabriz pour lequel on ne connaît pas d'activité}.
\end{blocgauche}
\end{frame}

%%%%%%%%%%%%%%%%%%%%
\begin{frame}[fragile]{L'opérateur \texttt{left outer join}}

\begin{redbullet}
   \item Renvoie tous les nuplets de la table directrice (à gauche)
   \item Associe à chaque nuplet un nuplet de la table de droite \alert{si un tel nuplet existe}
   \item Sinon, les attributs de la table de droite sont à \sqlkw{null}
\end{redbullet}

\vfill


\begin{minted}[baselinestretch=.9,
               fontsize=\small,
               framesep=2mm]{sql}
select nom, codeActivité
from Logement left outer join Activité 
       on (code = codeLogement)
\end{minted}

\vfill
On obtient, en plus de la jointure standard: 

\vfill
\begin{tabular}{c|c}
\hline
Tabriz & \sqlkw{null} \\ \hline 
\end{tabular}

\end{frame}



%%%%%%%%%%%%%%%%%%%%
\begin{frame}[fragile]{Le tri, \texttt{order by}}

On peut demander explicitement le tri du résultat sur
une ou plusieurs expressions avec la clause \texttt{order by}

\vfill


\begin{minted}[baselinestretch=.9,
               fontsize=\small,
               framesep=2mm]{sql}
   select *  
   from Logement 
   order by lieu, capacité
\end{minted}

\vfill

En ajoutant des clauses sur l'ordre du tri (\texttt{ascending} ou \texttt{descending})

\vfill

\begin{minted}[baselinestretch=.9,
               fontsize=\small,
               framesep=2mm]{sql}
   select *  
   from Logement 
   order by lieu desc, capacité desc
\end{minted}

\end{frame}
%%%%%%%%%%%
\begin{frame}{À retenir}

Les fondements du langage SQL (logique ou algèbre) sont une base sur laquelle
beaucoup d'extensions pratiques sont possibles.

\vfill

\begin{redbullet}
   \item Valeurs à \sqlkw{null}, jointures externes, tris
  \item  Mais aussi des fonctions, spécifiques à chaque système
  \item Ne change en rien \alert{l'interprétation} du langage, que vous devez maintenant maîtriser.
\end{redbullet}

\vfill
SQL est un standard, mais chaque système propose des spécificités au-delà du standard.

\end{frame}

